\section{Realización General del Proyecto}
    Primero se presentan las interpretaciones de los términos relacionados 
    a la detección de bordes. Estas interpretaciones se basan en conocimientos 
    previos y en lo visto en clases \parencite{tinoco2026pdi}: 
    \begin{itemize}
        \item $G_x = \frac{\partial I(x, y)}{\partial x}$: Indica cómo cambia la intensidad de los colores de forma horizontal. Estos cambios permiten encontrar líneas verticales.
        \item $G_y = \frac{\partial I(x, y)}{\partial y}$: Indica cómo cambia la intensidad de los colores de forma vertical. Estos cambios permiten encontrar líneas horizontales.
        \item $G = \sqrt{G_x^2+G_y^2}$: Representan la magnitud del cambio de la intensidad de los colores de forma direccional. Es la magnitud de la derivada direccional.
        \item $\theta = \arctan(\frac{G_y}{G_x})$: Es la dirección hacia donde la intensidad de los colores cambia ``más fuerte''. Es el ángulo de la derivada direccional.
    \end{itemize}

    La implementaciones de las funciones relacionadas a los filtros se basan en 
    las de CV2 \parencite{opencv_library}, es decir, se verificó que las funciones 
    programadas con NumPy \parencite{numpy_library} coincidan en valores con 
    las de CV2. Aunque estas funciones implementadas en NumPy nunca son llamadas, 
    debido a que son ``lentas'' en comparación con su equivalente en CV2 que se 
    ejecuta en Cpp. Por ello, solamente se manejan estas funciones con los parámetros 
    adecuados para generar los resultados esperados que son requeridos para la 
    práctica.\\ 

    Todas las matrices asociadas a los filtros revisados en clase fueron implementados 
    de forma manual, salvo los asociados a los suavizados gaussianos. Estos filtros 
    fueron comparados sobre las diferentes variantes de suavizados aplicadas a las 
    imágenes que, de forma general, permitieron reducir el ruido en regiones con 
    texturas de alto contraste que podrían confundirse con bordes, es decir, píxeles 
    aislados que son tratados como bordes. En algunos casos, el suavizado favoreció 
    a encontrar bordes más definidos, en otros simplemente los difuminaron.